\documentclass[11pt,a4paper]{article}

\usepackage[utf8]{inputenc}
\usepackage[T1]{fontenc}
\usepackage[a4paper,margin=2.5cm]{geometry}
\usepackage{array}
\usepackage{booktabs}
\usepackage{hyperref}

\title{HPC Tools Task 3: DGESV Profiling and Optimization}
\author{}
\date{July 2026}

\begin{document}

\maketitle

\section{Objective and Environment}

The goal of Task 3 was to benchmark and analyze the custom \texttt{dgesv}
implementation with Intel compilers, compare the results against the GCC
campaign from Task 2, profile the code, and identify optimizations that have
measurable impact.

The solver under study is the baseline Gaussian elimination with partial
pivoting plus back substitution implemented in \texttt{dgesv.c}. The reference
solver is LAPACKE \texttt{dgesv}; for the Intel campaign it is linked against
Intel MKL.

FT3 environment used for the Intel campaign:

\begin{itemize}
\item \texttt{intel/2021.3.0}
\item \texttt{imkl/2021.3.0}
\item \texttt{icc 2021.3.0}
\item \texttt{icx 2021.3.0}
\end{itemize}

The assignment requested \texttt{icx 2024.2.1}, but that version was not
visible in the \texttt{cesga/2025} module tree on 2026-07-05. We therefore ran
the experiments with the newest Intel toolchain that was actually available on
FT3 and documented that limitation explicitly.

\section{Methodology}

We followed the same general methodology as in Task 2:

\begin{itemize}
\item matrix sizes: \texttt{1024}, \texttt{2048}, \texttt{4096}
\item three runs per configuration
\item median execution time reported in the summary CSV
\item run-level data stored in a separate raw CSV
\end{itemize}

New Task 3 infrastructure:

\begin{itemize}
\item \texttt{run\_dgesv\_bench\_icc\_2021\_3\_0.sh}
\item \texttt{run\_dgesv\_bench\_icx.sh}
\item \texttt{profile\_dgesv\_intel.sh}
\end{itemize}

The Intel benchmark scripts also emit compiler vectorization reports.

For profiling and analysis we used:

\begin{itemize}
\item \texttt{perf stat -d}
\item \texttt{valgrind --leak-check=full}
\item \texttt{valgrind --tool=cachegrind}
\end{itemize}

VTune and Advisor modules are available on FT3, and the helper script supports
them, but the main conclusions below were already clear from \texttt{perf} and
Valgrind.

\section{Intel Benchmark Tables}

\subsection{ICC 2021.3.0}

\texttt{fast} deserves one important note. The literal \texttt{icc -fast}
driver path built a binary that crashed at runtime on FT3 with
\texttt{*** stack smashing detected ***}, even after fixing the MKL link path.
To keep the required \texttt{fast} table entry meaningful, we used the
fast-equivalent bundle:

\begin{center}
\texttt{-O3 -xHost -no-prec-div -fp-model fast=2 -ansi-alias}
\end{center}

The table label remains \texttt{fast}, but the report must mention this
FT3-specific compiler-driver defect.

\begin{table}[h]
\centering
\begin{tabular}{lrrrr}
\toprule
Matrix size & O2 & O3 & fast & Ref \\
\midrule
1024x1024 & 2145 & 1868 & 2027 & 118 \\
2048x2048 & 22533 & 22331 & 16745 & 418 \\
4096x4096 & 356421 & 368530 & 194918 & 2774 \\
\bottomrule
\end{tabular}
\end{table}

\subsection{ICX 2021.3.0}

\begin{table}[h]
\centering
\begin{tabular}{lrrrrr}
\toprule
Matrix size & O2 & O3 & Ofast & Ofast+ipo & Ref \\
\midrule
1024x1024 & 1916 & 1969 & 1915 & 2388 & 90 \\
2048x2048 & 22658 & 22082 & 21012 & 28892 & 481 \\
4096x4096 & 356030 & 369784 & 352261 & 452793 & 3113 \\
\bottomrule
\end{tabular}
\end{table}

\section{Combined View with Task 2}

The best GCC medians from Task 2 were:

\begin{itemize}
\item 1024: \texttt{2052 ms} with \texttt{GCC 8.4.0 Ofast-vec}
\item 2048: \texttt{22742 ms} with \texttt{GCC 8.4.0 Ofast-vec}
\item 4096: \texttt{400211 ms} with \texttt{GCC 8.4.0 Ofast-vec}
\end{itemize}

The best overall medians after combining Tasks 2 and 3 are:

\begin{table}[h]
\centering
\begin{tabular}{l l rr}
\toprule
Matrix size & Best overall configuration & Our solver (ms) & Reference (ms) \\
\midrule
1024x1024 & ICC 2021.3.0 O3 & 1868 & 88 \\
2048x2048 & ICC 2021.3.0 fast-equivalent & 16745 & 418 \\
4096x4096 & ICC 2021.3.0 fast-equivalent & 194918 & 2774 \\
\bottomrule
\end{tabular}
\end{table}

Relative to the best GCC result from Task 2, the Intel campaign improved the
custom solver by:

\begin{itemize}
\item $1.10\times$ at 1024 ($2052 / 1868$)
\item $1.36\times$ at 2048 ($22742 / 16745$)
\item $2.05\times$ at 4096 ($400211 / 194918$)
\end{itemize}

This means the Intel toolchain advantage becomes much larger as the problem
size grows.

\section{Performance Analysis}

\subsection{Main Findings}

The most important result is that the Intel compilers, especially ICC with the
fast-equivalent optimization bundle, outperform all GCC configurations from
Task 2 for the large matrices.

The clearest case is 4096x4096:

\begin{itemize}
\item best GCC result: \texttt{400211 ms}
\item best ICC result: \texttt{194918 ms}
\end{itemize}

That is a $2.05\times$ speedup without changing the algorithm, only the
compiler configuration and runtime library stack.

\subsection{Why the Intel Results Improve}

The main causes are:

\begin{itemize}
\item better autovectorization of the row-swap and elimination loops
\item more aggressive alias assumptions in the fast-equivalent ICC configuration
\item stronger MKL reference performance than the OpenBLAS-based reference from Task 2
\end{itemize}

The reference solver also benefits from the Intel stack. For example, at 4096
the best GCC reference median was \texttt{4694 ms} while ICC fast achieved
\texttt{2774 ms}, roughly $1.69\times$ faster.

\subsection{Most Unexpected Results}

Two results stood out as unexpected.

First, ICC fast was not the best configuration at 1024. ICC O3 achieved
\texttt{1868 ms}, while the fast-equivalent path took \texttt{2027 ms}. For
small matrices, the extra aggressiveness of the fast-equivalent flags does not
compensate for setup overhead and reduced cost-model selectivity.

Second, ICX Ofast+ipo was consistently worse than plain Ofast:

\begin{itemize}
\item \texttt{2388 ms} vs \texttt{1915 ms} at 1024
\item \texttt{28892 ms} vs \texttt{21012 ms} at 2048
\item \texttt{452793 ms} vs \texttt{352261 ms} at 4096
\end{itemize}

This suggests that IPO increased code size or hurt locality enough that the
extra whole-program optimization was not beneficial for this solver.

\subsection{Quantified Optimization Impact}

Within the Intel campaign, the best implemented optimization was the ICC
fast-equivalent flag bundle.

Compared with ICC O2, it improved the solver by:

\begin{itemize}
\item $1.06\times$ at 1024
\item $1.35\times$ at 2048
\item $1.83\times$ at 4096
\end{itemize}

Compared with ICC O3, it improved the solver by:

\begin{itemize}
\item $0.92\times$ at 1024, so slightly worse
\item $1.33\times$ at 2048
\item $1.89\times$ at 4096
\end{itemize}

The optimization is therefore clearly beneficial for medium and large problems,
which are the more relevant HPC cases.

\section{Vectorization Analysis}

\subsection{Are the Key Parts Autovectorized?}

Yes. The key numeric loops in \texttt{dgesv.c} are autovectorized by both Intel
compilers.

From the ICX Ofast report:

\begin{itemize}
\item row swap loop over \texttt{a} is vectorized
\item row swap loop over \texttt{b} is vectorized
\item elimination update over \texttt{a} is vectorized
\item elimination update over \texttt{b} is vectorized
\end{itemize}

From the Intel optimization report already present in the repository
(\texttt{ipo\_out.optrpt}), ICC also vectorizes:

\begin{itemize}
\item pivot search inner loop at \texttt{dgesv.c(8,5)}
\item row swap loops at \texttt{dgesv.c(17,7)} and \texttt{dgesv.c(23,7)}
\item elimination loops at \texttt{dgesv.c(34,7)} and \texttt{dgesv.c(38,7)}
\item one back-substitution inner loop at \texttt{dgesv.c(47,7)}
\end{itemize}

The ICC report even estimates a potential speedup of about $4.22\times$ for
the elimination update loops after vectorization.

\subsection{What Is Not Vectorized, and Why?}

Some loops still resist vectorization:

\begin{itemize}
\item \texttt{generate\_matrix()} is not vectorized because it calls \texttt{rand()}
\item \texttt{check\_result()} is not vectorized in the ICC fast report because it uses an unsigned induction variable and has early exit behavior
\item the outer pivot-search loop is not vectorized by ICX because the reduction and selected row index escape the loop
\item parts of back substitution remain limited by loop-carried dependences and by strided access to \texttt{b[j*nrhs + rhs]}
\end{itemize}

This behavior is consistent with the algorithm. The row-wise loops are friendly
to SIMD, but the dependency structure in pivoting and back substitution is
inherently harder.

\subsection{Can Source Changes Help?}

Yes, but the benefit is limited because the baseline already vectorizes the
main row-based loops under Intel compilers.

Possible source-level improvements are:

\begin{itemize}
\item replace unsigned loop counters in non-critical helper loops with signed integers
\item separate pivot search from row index updates more clearly to help ICX identify the reduction
\item use \texttt{restrict}-based helper routines similar to the separate autovectorization experiment from Task 2
\item improve alignment guarantees for matrix allocations
\end{itemize}

However, for the Task 3 baseline, compiler-flag tuning delivered larger gains
than additional source refactoring.

\section{Memory and Cache Behaviour}

\subsection{Heap Usage}

The first Valgrind pass showed that \texttt{main.c} leaked the five main
allocations:

\begin{itemize}
\item \texttt{a}
\item \texttt{b}
\item \texttt{aref}
\item \texttt{bref}
\item \texttt{ipiv}
\end{itemize}

We fixed that by freeing all owned buffers before returning from \texttt{main()}.

After the fix, Valgrind reports:

\begin{itemize}
\item definitely lost: 0 bytes
\item indirectly lost: 0 bytes
\item possibly lost: 0 bytes
\end{itemize}

There are still some reachable blocks inside MKL and loader internals, but the
program's own allocations are no longer leaked.

\subsection{Access Pattern}

The solver uses row-major matrices. That creates two very different memory
behaviors.

Good spatial locality appears in:

\begin{itemize}
\item row swaps over \texttt{a[i*n + j]}
\item row swaps over \texttt{b[i*nrhs + rhs]}
\item elimination updates over \texttt{a[k*n + j]}
\item elimination updates over \texttt{b[k*nrhs + rhs]}
\end{itemize}

These loops walk contiguous row segments and are exactly the loops that the
compilers vectorize most effectively.

Poorer locality appears in:

\begin{itemize}
\item pivot search, which reads \texttt{a[k*n + i]} and therefore traverses a column in row-major storage
\item back substitution, which accesses \texttt{b[j*nrhs + rhs]} and is also strided in the inner loop
\end{itemize}

These accesses have weaker spatial locality, create more cache pressure, and
are harder to vectorize efficiently.

\subsection{Cache Observations}

The profiling results support the idea that the solver becomes memory-sensitive
as the matrices grow.

\texttt{perf stat -d} on ICX Ofast at size 512 showed:

\begin{itemize}
\item IPC around 1.15
\item L1 data-cache miss rate around 33.5\%
\item LLC-load miss rate around 99.17\% of LLC loads
\end{itemize}

This is consistent with a kernel that streams large working sets and frequently
fetches data from deeper cache levels.

\texttt{cachegrind} on the optimized ICC fast-equivalent path at size 256 showed:

\begin{itemize}
\item D1 miss rate around 5.4\%
\item LL miss rate around 1.1\%
\end{itemize}

Those smaller-size results are qualitatively useful even if Valgrind is not
appropriate for absolute timing. They show that cache behavior is still
reasonable for smaller matrices, while the hardware \texttt{perf} counters
indicate much heavier pressure once the matrix grows.

\section{Profiling Summary}

The profiling tools used in practice were:

\begin{itemize}
\item \texttt{perf stat -d}
\item \texttt{valgrind --leak-check=full}
\item \texttt{valgrind --tool=cachegrind}
\end{itemize}

Main conclusions from profiling:

\begin{itemize}
\item the solver is not branch-miss dominated; branch-miss rates stayed low in the \texttt{perf} runs
\item the key bottleneck is memory traffic in the large row-update phases and the strided accesses in pivot search and back substitution
\item the custom solver allocates the correct matrix sizes, and after the memory fix it no longer leaks its own heap allocations
\end{itemize}

\section{Implemented Optimizations}

The main implemented optimizations in Task 3 were:

\begin{enumerate}
\item Intel benchmark infrastructure.

We added dedicated scripts for ICC and ICX, with median-of-three reporting and
vectorization-report emission.

\item ICC fast workaround.

Because the literal \texttt{icc -fast} path crashed at runtime on FT3, we
replaced it with the fast-equivalent flag bundle:

\begin{center}
\texttt{-O3 -xHost -no-prec-div -fp-model fast=2 -ansi-alias}
\end{center}

This preserved the intended aggressive optimization behavior while producing
stable binaries and the best large-matrix results.

\item Memory-management fix.

We added missing \texttt{free()} calls in \texttt{main.c}, eliminating the
definite leaks reported by Valgrind.

\item Reusable profiling helper.

We added \texttt{profile\_dgesv\_intel.sh} so that \texttt{perf},
\texttt{cachegrind}, \texttt{massif}, \texttt{advisor}, and \texttt{vtune}
can be run consistently on the Intel binaries.

\item Separate optimized experiment.

We created a separate experimental solver in \texttt{dgesv\_intel\_opt.c} so
that the baseline and the earlier helper-based variant remain untouched. This
version combines three ideas:

\begin{itemize}
\item \texttt{restrict}-qualified row pointers
\item manually unrolled contiguous row helpers for swaps, updates, and scaling
\item a rewritten back-substitution phase that operates on complete RHS rows instead of strided scalar updates
\end{itemize}

The complete single-run comparison under the same ICC fast-equivalent flags
used in the main Task 3 campaign is shown in Table~\ref{tab:separate-opt}.

\begin{table}[h]
\centering
\begin{tabular}{rrrr}
\hline
Size & Older helper-based variant (ms) & \texttt{dgesv\_intel\_opt.c} (ms) & Speedup \\
\hline
512 & 176 & 81 & 2.17x \\
1024 & 1731 & 641 & 2.70x \\
2048 & 16234 & 6663 & 2.44x \\
4096 & 194615 & 62337 & 3.12x \\
\hline
\end{tabular}
\caption{Separate optimized experiment under ICC 2021.3.0 fast-equivalent flags.}
\label{tab:separate-opt}
\end{table}

The dedicated artifact for this comparison is stored in
\texttt{dgesv\_results\_icc\_separate\_opt.csv}.

These results strengthen the earlier conclusion: back substitution was a larger
bottleneck than the row-update kernels alone, and reorganizing the solve to
operate on contiguous RHS rows is substantially more effective than only adding
helper-based swap and update routines. The benefit also scales well with
problem size, exceeding 3x at 4096.
\end{enumerate}

\section{Final Conclusion}

The Task 3 results show that compiler and runtime choice matter a lot for this
solver. The Intel toolchain produced the best results of the whole assignment
campaign, especially for the larger matrices where the gap over GCC became
large.

The best overall configuration in this repository is:

\begin{itemize}
\item ICC 2021.3.0 fast-equivalent for 2048 and 4096
\item ICC 2021.3.0 O3 for 1024
\end{itemize}

The main reasons are:

\begin{itemize}
\item successful autovectorization of the row-based hot loops
\item aggressive but still stable floating-point and aliasing optimizations in the ICC fast-equivalent path
\item stronger MKL performance for the reference solver
\end{itemize}

The code is still fundamentally limited by the algorithm's memory behavior:
column-like pivot search and strided back substitution remain difficult for
caches and SIMD. That is why profiling points to memory pressure rather than
control-flow overhead as the main scaling problem.

If more time were available, the next experiments we would prioritize are:

\begin{itemize}
\item running the same campaign on a true \texttt{icx 2024.2.1} installation
\item benchmarking the \texttt{dgesv\_autovec.c} helper-based variant under Intel compilers
\item testing aligned allocation and explicit SIMD hints on the pivot and back-substitution kernels
\end{itemize}

\end{document}